\begin{papercard}{cao2024}
\cf{Stages} \stagetag{H} (which data practices of a cycle-tracking app women find most concerning, and why), \stagetag{C} (a mixed qualitative codebook with two coders), \stagetag{I} (reliability reported as ``Cohen's $\kappa$ of 87.5\,\%'', see limits), \stagetag{S}. A CHI 2024 vignette survey of 183 US women after \emph{Dobbs} (the overturn of \emph{Roe v.\ Wade}, June 2022): four manipulated data-practice factors, an ordinal mixed-model analysis and a qualitative content analysis of open answers. It is the closest user-side study of the thesis' app domain and models how to measure which features of a data practice drive concern.

\cf{Motivation} Period-tracking apps are the most popular women's mHealth apps and collect data (cycle, pregnancy, sex life, location) from which abortion can be inferred, and the data are widely shared with third parties; after \emph{Dobbs}, ``law enforcement can now request fertility-related records from period-tracking app companies as evidence of crimes'' \pc{cao2024}{1}{law enforcement can now request fertility-related records from period-tracking app companies as evidence of crimes}. Background cited: a scoping review of 23 women's mHealth apps (all permitted behavioural tracking, 61\,\% location tracking, 87\,\% shared data with third parties; 30\,\% showed no privacy policy) \pc{cao2024}{3}{In a scoping review of 23 popular women’s mHealth apps, it was found that all apps permitted behavioral tracking, while 61\% allowed location tracking}; ORCHA: 67\,\% of tested period apps share data for ``legal obligations''; the FTC case against Flo Health for sharing sensitive data with Facebook, Google, AppsFlyer and Flurry while telling users it did not \pc{cao2024}{3}{accusing them of allegedly sharing users’ sensitive health data with Facebook, Google, AppsFlyer, and Flurry over an extended period of time while they explicitly told their users they did not}; EFF's caution that bought datasets are not currently the main route to prosecution (police rely on messages and searches) \pc{cao2024}{3}{The Electronic Frontier Foundation (EFF) argues that although anyone may buy certain period-tracking apps’ datasets, it is not the primary strategy being used to criminalize abortion seekers}. Gap: earlier work outside the US or small interview studies; no feature-level quantification.

\cf{Research questions and design} RQ1 factors influencing privacy perceptions and practices; RQ2 knowledge of the overturn and its impact on concern; RQ3 expectations of privacy features and of stakeholders \pc{cao2024}{2}{RQ3: What are women’s expectations for privacy-enhancing features, actions, and information from stakeholders}. Design: fractional-factorial vignette study, where each participant answers about a subset of scenarios \pc{cao2024}{4}{In a fractional- factorial vignette study, each participant only answers questions about a subset of scenarios}. \emph{Four factors}: Collected Data (menstrual cycle data alone; plus precise location; plus mental-health data; plus physical-health data; plus intimacy data), Data Storage (device or cloud), Data Sharing (none; app company; app company plus law enforcement; plus third-party advertisers; plus healthcare providers), User Control (none; opt out of sharing; delete data). Levels were chosen from prior literature and ordered from low to high protection. 150 combinations minus unrealistic ones (e.g., ``not shared'' together with ``opt out from sharing'') = \textbf{75 scenarios}; each participant saw four random ones \pc{cao2024}{5}{Each participant was presented with four randomly selected scenarios}, in fixed order of factors, followed by a concern rating, a reason, and one attention-check item per scenario; participants missing more than two checks were dropped \pc{cao2024}{5}{we removed participants who missed more than two attention-check questions}. Further survey blocks: period-tracking use, concerns and mitigation, knowledge of app practices, familiarity with the overturn and its perceived impact, demographics. Pilot with 20 Prolific participants shortened the survey from six vignettes (31 min) to four (14 min) \pc{cao2024}{6}{we reduced the number of presented vignettes from six scenarios (average completion time of 31 minutes) to four scenarios}.

\cf{Sample} Prolific, US women (including trans women), 18+, approval rate $>$95\,\%; half from 15 abortion-ban states and half from 8 states with abortion allowed without gestational limit \pc{cao2024}{6}{we recruited half of our participants from 15 abortion-banned states (e.g., Texas), while the other half were recruited from 8 abortion-allowed states}; \$4 each; 200 recruited, 17 removed (failed checks, incomplete consent, mixed-policy state), \textbf{183} analysed \pc{cao2024}{7}{As a result, we removed 17 responses}; mean age 39; 51\,\% in legal states, 48\,\% in ban states; under-representation of Black (4.9\,\% vs 13.6\,\%) and Hispanic (2.7\,\% vs 19.1\,\%) women relative to census.

\cf{Analysis} \emph{Quantitative}: ordinal regression by a cumulative link mixed model (CLMM) with a random intercept per participant because of four repeated ratings; forward model selection by AIC among the four factors, scenario order, demographics and first-order interactions \pc{cao2024}{6}{we conducted an ordinal regression analysis by constructing a cumulative link mixed model (CLMM)}. \emph{Qualitative}: content analysis of open answers; one author coded 20 responses (10\,\%) to build the codebook, two researchers then coded independently and revised it over several meetings; reported Cohen's $\kappa$ ``87.5\,\%'' (``almost perfect''); the final codebook has 10 main codes, 55 sub-codes, 101 sub-sub-codes and 9 further sub-codes \pc{cao2024}{6}{The resulting codebook contains 10 main codes, 55 sub-codes, 101 sub-sub-codes, and 9 sub-sub-sub-codes}. IRB approval; positionality statement (four of five authors cisgender women, three based in the US).

\cf{Results (quantitative, Table 3)} Factor importance by AIC: Data Sharing most important, User Control second, Collected Data third, Data Storage least \pc{cao2024}{8}{the party with whom the data is being shared with (Data Sharing) was the most important factor}. Odds ratios of higher concern (baseline in brackets): sharing with law enforcement 93.3 (none), third parties 26.8, healthcare providers 7.4, app company 4.5; collected data: location 5.6, intimacy 5.2, mental health 2.3 (cycle data only); physical health 1.5 n.s.; user control: delete 0.08, opt-out of sharing 0.29 (none); cloud versus device 0.83 n.s.; Republican 3.0 and bachelor's degree 3.9 relative to baselines; interaction: deletion control reduces concern less when data go to healthcare providers or law enforcement \pc{cao2024}{9}{Table 3: CLMM regression model to describe how various scenario factors impact participants’ level of privacy concerns toward period-tracking data collection and use}. Device type for tracking made no difference \pc{cao2024}{8}{Using period-tracking apps either on a wearable device, a personal computer/tablet, or a phone generated similar levels of concern}; posting about fertility topics online worried them most, searching online or messaging least, even though law enforcement currently relies on these.

\cf{Results (descriptive and qualitative)} 67\,\% (123/183) track or tracked periods, 51\,\% use apps (Flo, Period Tracker, Apple, Clue) \pc{cao2024}{7}{67\% of our survey participants (123/183) reported that they are currently tracking their periods or have previously tracked their periods}; 74\,\% would accept an app recommended by healthcare professionals, most would not follow recommendations from government, employers or insurers \pc{cao2024}{8}{74\% of our participants (136/183) were at least somewhat likely to accept the app recommendations from healthcare professionals}. Sharing with law enforcement ``unacceptable'' for 37\,\% (67/183); third-party access worried 31\,\% (insurance, advertisers); 11\,\% (21/183) distrusted opt-out promises \pc{cao2024}{10}{11\% of participants (21/183) explicitly said that app companies might still share their data despite opt-out}. Only 8\,\% (14/183) attributed concern to the overturn when not prompted \pc{cao2024}{8}{only 8\% of participants (14/183) attributed their privacy concerns to the overturn of Roe v. Wade}; once asked, 60\,\% (110/183) said it increased their concern, 38\,\% saw minimum or no impact; 31\,\% (56/183) had switched or stopped apps, mostly for usability (54\,\%) then privacy (20\,\% of reasons; 5 of 183 because of the overturn). Mitigation: only 9\,\% (16/183) took steps (6 deleted the app) and 91\,\% (167/183) none, mainly lacking knowledge \pc{cao2024}{11}{91\% of our participants (167/183), however, mentioned that they had never used any mitigation strategies, mainly due to their lack of privacy knowledge and awareness}. Responsibility: 87\,\% named app developers, 49\,\% users, 48\,\% app-store companies, 18\,\% government \pc{cao2024}{12}{87\% of participants (159/183) considered app developers to be most responsible for protecting their period- tracking apps’ data privacy}. Wishes (n=62): usable controls (26\,\%), end-to-end encryption (19\,\%), granular permissions (15\,\%), anonymisation (6\,\%); transparency about purposes, with whom data are shared, benefits and risks, available controls, recoverability after deletion \pc{cao2024}{13}{let users have more control over data deletion and sharing}.

\cf{Discussion and recommendations} Concern is driven mostly by who receives the data; collection or sharing of data ``unnecessary'' for tracking (location, mental health, intimacy) is rejected. The \emph{My Body, My Data Act} limits sale/sharing of reproductive data ``except as is strictly necessary to provide a product or service'' without defining necessity or addressing law-enforcement requests \pc{cao2024}{14}{except as is strictly necessary to provide a product or service}. Recommendations: transparency mechanisms (policies, nutrition labels) that say how law-enforcement requests are handled; granular and truly usable controls, avoiding the ``privacy communication game'' \pc{cao2024}{14}{avoid the pitfalls of the privacy communication game}; pseudonyms or no account; policies that state whether law enforcement can request access \pc{cao2024}{15}{policies should clearly specify if law enforcement can request access to reproductive health data from companies}; public and school education.

\cf{Limits (authors and card)} Vignettes with a limited factor set and not real behaviour; Prolific bias and under-representation of Black and Hispanic women; only female participants (partner surveillance not studied) \pc{cao2024}{7}{Hispanic or Latino only constituted 2.7\% of our sample}. Card additions: the reliability figure is reported as ``Cohen's $\kappa$ of 87.5\,\%'', which looks like a percent-agreement value labelled $\kappa$ and is applied across a four-level codebook, so no per-code or chance-corrected statistic can be reconstructed; coding was by two researchers after one author drafted the codebook; ``unnecessary'' is a participant label, not a measured construct; factors confound ``who gets data'' with severity because the baseline is ``not shared'', so odds ratios are large and depend on that baseline; US-law context.

\cf{Use for this thesis} (1) Domain evidence that recipient (law enforcement, advertisers, insurers) and extra data types (location, intimacy, mental health) drive concern more than storage location, which supports ranking harm-tree branches by recipient and data type and checking apps against these combinations. (2) A concrete vignette factor set (data collected, storage, recipient, control) that can be adapted to code or compare real app policies, or later to a small user study; the thesis' CI-style coding of policy flows can reuse the same four dimensions. (3) A caution for reporting inter-coder reliability: state the statistic, the unit and the categories; the thesis should report a chance-corrected value per code. (4) Shows the mismatch between concerns and mitigation (91\,\% no steps), relevant for recommendations. (5) Use together with Moniz and Almeida (design probe), Malki et al.\ (post-Roe app practices) and Mehrnezhad and Almeida.
\end{papercard}
